% ===============================================================
% Section VIII  --  Conclusion
% Owner: Integration and Reporting Lead
% ===============================================================
\section{Conclusion}
\label{sec:conclusion}

This project set out to adapt a small, locally-hosted language model to a
bounded domain, three Shakespeare plays, for a user with no prior
exposure to the texts. We met the assignment's core objectives: a
structured corpus prepared at two granularities and enriched with
modern-English metadata; a mandatory RAG pipeline in which retrieved
evidence is demonstrably used and always displayed; four interaction
modes including a clearly-labelled stylised mode; a baseline-versus-RAG
comparison; and a structured evaluation with explicit failure analysis.
The system runs on a standard laptop with cached indices and no GPU,
satisfying the compute and deployability constraints.

The evaluation revealed a clean separation between the two halves of the
pipeline. Retrieval is the system's strength: across both scene- and
utterance-level indices the right passages were consistently surfaced,
and because retrieval is independent of the generation model, this result
carries over to the deployed Gemma-based system. Generation was the
limiting factor in the batch evaluation, where a lightweight stand-in
generator failed to ground its answers in evidence it had been given.

The most important design choices were therefore about retrieval and
honesty of comparison: enriching short archaic utterances with
modern-English context so they embed meaningfully, keeping the baseline
prompt-only so the contribution of retrieval was measurable, and always
showing evidence so a beginner could verify every claim. The clearest
lesson is that retrieval quality and generation quality must be evaluated
and improved separately; strong retrieval is wasted on a weak generator.

With more time, the priority is to re-run the evaluation harness against
the deployed Gemma~3~4B model so the strong retrieval scores are matched
by coherent answers, and to replace the scene-level diversity cap with a
minimum-coverage rule that guarantees act-level breadth without discarding
the most relevant passages on focused questions. A systematic sweep over
$k$ and a comparison of embedding models would further strengthen the
evaluation. None of these change the system's architecture; they tighten
the alignment between what the system demonstrably does and what the
evaluation measures, which is the standard this assignment rightly sets.
